\chapter{染色体}\label{ch:g9:chromosomes}

\Cref{rem:g9:heredity-and-traits:question}已把
决定因素锁定在\omterm{def:g6:cell-unit-of-life:parts}{细胞核}中。在恰当的时刻观察\omterm{def:g5:first-look-at-cells:cell}{细胞}，
\omterm{def:g6:cell-unit-of-life:parts}{细胞核}便会“露出底牌”：其内容物卷缩成可数、可拍照的
小体——\omterm{def:g9:chromosomes:chromosomes}{染色体}。
数一数体\omterm{def:g5:first-look-at-cells:cell}{细胞}和\omterm{def:g8:sexual-reproduction:gametes}{配子}中的\omterm{def:g9:chromosomes:chromosomes}{染色体}，便能解释家族相册里
几乎所有的疑问：\omterm{def:g8:sexual-reproduction:gametes}{配子}为何减半、\omterm{def:g9:chromosomes:chromosomes}{染色体}为何成对、
后代为何会重新组合，
以及决定婴儿是男孩还是女孩的那次“掷硬币”。

\section{捕捉细胞核的瞬间}

\begin{definition}[染色体]\label{def:g9:chromosomes:chromosomes}
\emph{染色体}\index{染色体}是\omterm{def:g6:cell-unit-of-life:parts}{细胞核}中的线状小体，
分裂时可见（经良好染色后），
此时内容物卷缩以便运输。两次分裂之间，
同样的丝状物松散舒展、不可见——工作态；
分裂时则凝缩——运输态。
经拍照并分类后，一个\omterm{def:g5:first-look-at-cells:cell}{细胞}的染色体组成其
\emph{核型}\index{核型}：
整副染色体依次摆开。
\end{definition}

\begin{proposition}[人体染色体数目]\label{prop:g9:chromosomes:counts}
\omterm{def:g9:chromosomes:chromosomes}{核型}给出三个数字，
贯穿全章：
\begin{enumerate}
  \item 人体每个体\omterm{def:g5:first-look-at-cells:cell}{细胞}都含有\emph{46}条\omterm{def:g9:chromosomes:chromosomes}{染色体}
——\omterm{def:g1:the-five-senses:sense}{皮肤}、\omterm{prop:g7:digestion-and-nutrients:absorption}{肝脏}、\omterm{def:g8:nervous-communication:neuron}{神经元}\omterm{def:g5:first-look-at-cells:cell}{细胞}皆如此
        （\cref{prop:g5:first-look-at-cells:growth}的
        分裂将整副\omterm{def:g9:chromosomes:chromosomes}{染色体}忠实地复制到每个\omterm{def:g5:first-look-at-cells:cell}{细胞}中）；
  \item 摆开后，46 条\omterm{def:g9:chromosomes:chromosomes}{染色体}分为\emph{23 对}，两条
        同源\omterm{def:g9:chromosomes:chromosomes}{染色体}大小相同、
        带型一致；
  \item \emph{\omterm{def:g8:sexual-reproduction:gametes}{配子}}打破此规律：\omterm{def:g4:animal-reproduction:sexual}{卵细胞}与
        \omterm{def:g8:reproductive-systems:male}{精子}各含\emph{23}条——每对中各取一条，
        每对均不完整。
\end{enumerate}
\end{proposition}
\admitted

\begin{example}[受精的算术]\label{ex:g9:chromosomes:arithmetic}
数目环环相扣，恰如一道证明。\omterm{def:g8:sexual-reproduction:gametes}{配子}各含 23 条；
\omterm{def:g5:flower-to-fruit:steps}{受精}（\cref{prop:g8:sexual-reproduction:core}）
两者相加：\omterm{prop:g8:sexual-reproduction:core}{受精卵}含 $23 + 23 = 46$——\omterm{def:g7:organs-at-work:recovery}{恢复}成对，
每对中一条来自母亲，一条
来自父亲。此后，孩子的每个体\omterm{def:g5:first-look-at-cells:cell}{细胞}都携带
父母双方的贡献——这正是
\cref{rem:g4:animal-reproduction:mix}从外部观察到的
现象：小马驹与其白斑，前面章节已见过。
\end{example}

\begin{remark}[配子为何必须减半]\label{rem:g9:chromosomes:whyhalve}
假设不然：含 46 条\omterm{def:g9:chromosomes:chromosomes}{染色体}的\omterm{def:g8:sexual-reproduction:gametes}{配子}将产生
含 92 条\omterm{def:g9:chromosomes:chromosomes}{染色体}的子代，184 条\omterm{def:g9:chromosomes:chromosomes}{染色体}的孙代——
\omterm{def:g9:chromosomes:chromosomes}{染色体}数目逐代翻倍。\omterm{def:g8:sexual-reproduction:gametes}{配子}特有的
减数分裂——使每个\omterm{def:g8:sexual-reproduction:gametes}{配子}各得\emph{每对中的一条
}，正是它使\omterm{def:g5:biodiversity:species}{物种}数目稳定在
46 条，代代如此。另需注意一项额外的收获：
\omterm{def:g8:sexual-reproduction:gametes}{配子}从每对中获得\emph{哪一条}，逐对独立决定。
一个亲本可产生 $2^{23}$ 种可能的\omterm{def:g9:chromosomes:chromosomes}{染色体}组合——
这正是
\cref{prop:g8:sexual-reproduction:shuffle}所述重组的
最终计数。
\end{remark}

\section{染色体携带决定因素}

\begin{proposition}[染色体确有作用的证据]\label{prop:g9:chromosomes:evidence}
\omterm{def:g9:chromosomes:chromosomes}{染色体}\emph{确是}决定因素的载体，这并非假设，
而是有证据支持：
\begin{itemize}
  \item 它们的行为恰好符合
        \cref{prop:g9:heredity-and-traits:conclusions}的要求：
        即：存在于每个\omterm{def:g5:first-look-at-cells:cell}{细胞}中，减半后进入
        \omterm{def:g8:sexual-reproduction:gametes}{配子}，\omterm{def:g5:flower-to-fruit:steps}{受精}时复原，
        可按同源\omterm{def:g9:chromosomes:chromosomes}{染色体}逐对选择；
  \item \omterm{def:g9:chromosomes:chromosomes}{核型}异常证明了其“载荷”：\omterm{def:g9:chromosomes:chromosomes}{染色体}数目
        多一条或少一条，便会改变整个
        身体的发育——比如较小的第 21 对多出一份，
        便产生可辨认的一组\omterm{def:g9:heredity-and-traits:traits}{性状}，
        即唐氏综合征。一条\omterm{def:g9:chromosomes:chromosomes}{染色体}数目出错，
        众多\omterm{def:g9:heredity-and-traits:traits}{性状}便随之改变：这些丝状物
        携带着\omterm{def:g9:heredity-and-traits:traits}{性状}指令；
  \item 另有一对\omterm{def:g9:chromosomes:chromosomes}{染色体}在每份\omterm{def:g9:chromosomes:chromosomes}{核型}中都“露面”，
        即\emph{性\omterm{def:g9:chromosomes:chromosomes}{染色体}}。
\end{itemize}
\end{proposition}
\admitted

\begin{definition}[X、Y 与抛硬币]\label{def:g9:chromosomes:xy}
第 23 对有两种形式：两条相同的\emph{X}
\omterm{def:g9:chromosomes:chromosomes}{染色体}——女孩的\omterm{def:g9:chromosomes:chromosomes}{核型}；或一条 X 加一条
较小的\emph{Y}——男孩的\omterm{def:g9:chromosomes:chromosomes}{核型}。减半之后，便是一次“抛硬币”：
每个\omterm{def:g4:animal-reproduction:sexual}{卵细胞}都带有一条 X；\omterm{def:g8:reproductive-systems:male}{精子}则携带 X 或 Y，
各占一半。X \omterm{def:g8:reproductive-systems:male}{精子}使\omterm{prop:g8:sexual-reproduction:core}{受精卵}成为 XX，Y \omterm{def:g8:reproductive-systems:male}{精子}
使其成为 XY——机会各为二分之一
（\cref{ch:g9:heredity-and-traits}中的机会语言，
终于有了其机制），由哪个\omterm{def:g8:reproductive-systems:male}{精子}先抵达
卵子决定，发生在\omterm{def:g5:flower-to-fruit:steps}{受精}的瞬间。
\end{definition}

\begin{omfigure}
\begin{tikzpicture}[scale=0.9]
  % mother
  \node[draw, thick, omDef, rounded corners=4pt, font=\small,
    align=center] (mo) at (1.2,4.2) {母亲 \\ XX};
  \node[draw, thick, omDef, rounded corners=4pt, font=\small,
    align=center] (fa) at (7.8,4.2) {父亲 \\ XY};
  % gametes
  \node[draw, omDef, circle, font=\footnotesize] (e1) at (1.2,2.4) {X};
  \node[draw, omDef, circle, font=\footnotesize] (s1) at (6.6,2.4) {X};
  \node[draw, omDef, circle, font=\footnotesize] (s2) at (9.0,2.4) {Y};
  \draw[->, thick] (mo) -- (e1);
  \draw[->, thick] (fa) -- (s1);
  \draw[->, thick] (fa) -- (s2);
  \node[font=\footnotesize, left] at (0.8,2.4) {每个卵细胞};
  \node[font=\footnotesize, right, align=left] at (9.4,2.4)
    {精子：\\ 一半 X，一半 Y};
  % offspring
  \node[draw, thick, omThm, rounded corners=4pt, font=\small,
    align=center] (girl) at (3.4,0.4) {XX\\ 女孩};
  \node[draw, thick, omMeth, rounded corners=4pt, font=\small,
    align=center] (boy) at (5.8,0.4) {XY\\ 男孩};
  \draw[->, thick, omExo] (e1) -- (girl);
  \draw[->, thick, omExo] (s1) -- (girl);
  \draw[->, thick, omExo] (e1) -- (boy);
  \draw[->, thick, omExo] (s2) -- (boy);
  \node[font=\footnotesize, align=center] at (4.6,-0.5)
    {机会各半，发生于受精时};
\end{tikzpicture}

{\small 第 23 对的掷硬币：卵子的 X 与一个含 X 或含 Y 的\omterm{def:g8:reproductive-systems:male}{精子}相遇——
\omterm{def:g9:chromosomes:chromosomes}{核型}便随这次相遇而定。}
\end{omfigure}

\begin{example}[旧问题，短答案]\label{ex:g9:chromosomes:answers}
这副牌解决了历史旧账。孩子是男是女，
与父母任何一方的意愿、饮食或季节都无关，
取决于那次掷币，
其机制见\cref{def:g9:chromosomes:xy}（至于旧时的不公：
决定性的差异落在
\emph{\omterm{def:g8:reproductive-systems:male}{精子}}上）。两方各出一半，也回答了家庭相册里的问题：
母亲与父亲各提供一副完整的半副牌，
各为 23 条，\omterm{def:g9:heredity-and-traits:heredity}{遗传}的算术
对两条线一视同仁。
\end{example}

\begin{method}[读懂核型图]\label{met:g9:chromosomes:read}
拿到一张\omterm{def:g9:chromosomes:chromosomes}{核型}照片时：
\begin{enumerate}
  \item 计数：体\omterm{def:g5:first-look-at-cells:cell}{细胞}中应为 46，\omterm{def:g8:sexual-reproduction:gametes}{配子}中应为
        23；
  \item 配对：依大小与带型配对同源\omterm{def:g9:chromosomes:chromosomes}{染色体}，
        共 22 对常\omterm{def:g9:chromosomes:chromosomes}{染色体}加上第 23 对；
  \item 读第 23 对：XX 或 XY；
  \item 核对数目：凡成员为一条或三条的对，
        依
        \cref{prop:g9:chromosomes:evidence}，
        说明发育出现了偏移。
\end{enumerate}
\end{method}

\section{练习}

\begin{exercise}[$\star$]\label{exo:g9:chromosomes:1}
什么是\omterm{def:g9:chromosomes:chromosomes}{染色体}？它何时可见？
\end{exercise}

\begin{exercise}[$\star$]\label{exo:g9:chromosomes:2}
给出人类的三个数目，并说明各自成立的位置。
\end{exercise}

\begin{exercise}[$\star$]\label{exo:g9:chromosomes:3}
算一算\omterm{def:g5:flower-to-fruit:steps}{受精}过程中的数目变化，
并说明它\omterm{def:g7:organs-at-work:recovery}{恢复}了什么。
\end{exercise}

\begin{exercise}[$\star$]\label{exo:g9:chromosomes:4}
为什么\omterm{def:g8:sexual-reproduction:gametes}{配子}必须把牌数减半？
若不然，会出现怎样的失控？
\end{exercise}

\begin{exercise}[$\star$]\label{exo:g9:chromosomes:5}
第 23 对的两种版本是什么？哪一方的
\omterm{def:g8:sexual-reproduction:gametes}{配子}决定二者之间的取舍？
\end{exercise}

\begin{exercise}[$\star$]\label{exo:g9:chromosomes:6}
多出一条 21 号\omterm{def:g9:chromosomes:chromosomes}{染色体}会产生什么，
这说明了关于“线”的什么？
\end{exercise}

\begin{exercise}[$\star\star$]\label{exo:g9:chromosomes:7}
\cref{prop:g8:sexual-reproduction:shuffle}所述的洗牌究竟
发生在本章所说的哪一环节，
又一方的洗牌空间
有多大？
\end{exercise}

\begin{exercise}[$\star\star$]\label{exo:g9:chromosomes:8}
将\cref{prop:g9:heredity-and-traits:conclusions}的每一条要求，
与满足它的
\omterm{def:g9:chromosomes:chromosomes}{染色体}行为一一配对。
\end{exercise}

\begin{exercise}[$\star\star$]\label{exo:g9:chromosomes:9}
“每次出生为男孩的机会都是二分之一，
无论之前出现过什么。”请以其机制为之辩护——
为什么连续三个女孩并不改变任何东西？
\end{exercise}

\begin{exercise}[$\star\star$]\label{exo:g9:chromosomes:10}
对以下情形运用\cref{met:g9:chromosomes:read}：(a) 46,XX 的体\omterm{def:g5:first-look-at-cells:cell}{细胞}；
(b) 23,Y 的\omterm{def:g8:sexual-reproduction:gametes}{配子}；(c) 体\omterm{def:g5:first-look-at-cells:cell}{细胞}为 47 条，
多出的一条在 21 号对中。
\end{exercise}

\begin{exercise}[$\star\star$]\label{exo:g9:chromosomes:11}
历史上总把王国生女归咎于王后。
请借助\cref{ex:g9:chromosomes:answers}，
逐个\omterm{def:g8:sexual-reproduction:gametes}{配子}地纠正这一记载。
\end{exercise}

\begin{exercise}[$\star\star\star$]\label{exo:g9:chromosomes:12}
身体的每个\omterm{def:g5:first-look-at-cells:cell}{细胞}都带有完整的 46 条——
然而\omterm{def:g8:nervous-communication:neuron}{神经元}与\omterm{def:g1:the-five-senses:sense}{皮肤}\omterm{def:g5:first-look-at-cells:cell}{细胞}却大不相同
（\cref{rem:g6:cell-unit-of-life:twoways}）。请提出由此产生的
疑问，并说明其解决的方向
（这是下一章的内容：\omterm{def:g9:chromosomes:chromosomes}{染色体}上的什么是
\emph{被读取}的）。
\end{exercise}

\section{问题：细胞遗传学实验台}

\begin{problem}[{周末问题——同一张灯箱台上的四份
核型}]\label{pb:g9:chromosomes:1}
某医院\omterm{def:g5:first-look-at-cells:cell}{细胞}\omterm{def:g9:heredity-and-traits:heredity}{遗传}学实验室的教学台：\omterm{def:g9:chromosomes:chromosomes}{核型} A
为 46 条\omterm{def:g9:chromosomes:chromosomes}{染色体}，第 23 对为 XY；B 为 46 条，第 23 对为 XX；
C 为 23 条\omterm{def:g9:chromosomes:chromosomes}{染色体}，其中有一条 X；D 为 47 条，
第 21 对有三份拷贝，第 23 对为 XX。

\medskip\noindent\textbf{第一部分——读图。}
\begin{enumerate}
  \item 按
        \cref{met:g9:chromosomes:read}读取 A、B 和 C，
        判断它们各可能取自谁的\omterm{def:g5:first-look-at-cells:cell}{细胞}？
  \item C 的数目不是误差。它是哪种\omterm{def:g5:first-look-at-cells:cell}{细胞}，
        其 23 条\omterm{def:g9:chromosomes:chromosomes}{染色体}记录了怎样的
        分裂历史？
  \item 读 D：数目、配对，以及它所预示的
        发育情况。
  \item 哪一方亲本的\omterm{def:g8:sexual-reproduction:gametes}{配子}决定了 A 的第 23 对\omterm{def:g9:chromosomes:chromosomes}{染色体}，
        在那次\omterm{def:g5:flower-to-fruit:steps}{受精}中，这种结果的概率是多少？
\end{enumerate}

\medskip\noindent\textbf{第二部分——计算
问题。}
\begin{enumerate}[resume]
  \item 一名学员问，为什么同一人的体\omterm{def:g5:first-look-at-cells:cell}{细胞}
        都与\omterm{def:g9:chromosomes:chromosomes}{核型} A 一致？请借助
        复制性分裂来作答，见
        \cref{prop:g5:first-look-at-cells:growth}。
  \item 另一名学员问，两个各有 46 条\omterm{def:g9:chromosomes:chromosomes}{染色体}的亲本
        怎样生出了\omterm{def:g9:chromosomes:chromosomes}{染色体}只有 23 条的 C？请说出那种特殊分裂及其规律。
  \item 将 C 与一个携带 X 的\omterm{def:g4:animal-reproduction:sexual}{卵细胞}结合，再与
        本实验台上携带 Y 的邻居样本结合——等等：C
        \emph{本身不含} Y。请说明 C 的 X
        能为其参与产生的任何孩子保证什么。
  \item 计算一个亲本所能产生的
        全部可能\omterm{def:g8:sexual-reproduction:gametes}{配子}：说明逐对独立，
        并得出 $2^{23}$ 的结论——再考虑
        \omterm{def:g5:flower-to-fruit:steps}{受精}时两个这种\omterm{def:g8:sexual-reproduction:gametes}{配子}的配对。
\end{enumerate}

\medskip\noindent\textbf{第三部分——咨询
室。}
\begin{enumerate}[resume]
  \item \omterm{def:g9:chromosomes:chromosomes}{核型}为 D 的孩子的父母问：“我们
        做错了什么？”请就减数分裂的意外给出
        坦率的回答，并说明这个孩子的 47 条\omterm{def:g9:chromosomes:chromosomes}{染色体}
        究竟改变了什么、又没有改变什么，
        ——对于屋里的爱。（以咨询师身份作答，
        生物学内容须准确。）
  \item 一对已有三个女儿的夫妇问，下一个孩子是儿子的
        概率是多少？请用掷硬币作答，并纠正
        “该轮到男孩了”的直觉。
  \item 一位参观者声称，饮食可以选择婴儿的性别。
        请依据
        \cref{def:g9:chromosomes:xy}审查这一说法：
        性别由何处决定？饮食又能影响什么？
  \item 用一句话作结：这四张照片合起来
        证明了
        \omterm{def:g9:heredity-and-traits:heredity}{遗传}究竟依托于何处。
\end{enumerate}
\end{problem}
